% DOC NO. RL-Z0-SELF-HOSTING · REV. A
%
% This file IS the document. Edit it directly - ripdoc compiles it
% as it stands and stamps the provenance line at build time.
%
%   ripdoc build --only docs-self-hosting
%
% Promoted from docs/self-hosting.md on 2026-09-07 by `ripdoc promote`.
\documentclass[fontpath=fonts/,logopath=logo/]{riposte-doc}
\usepackage{riposte-md}

\title{CHANNEL Z0 — Self-Hosting on a Homelab (Proxmox \& TrueNAS)}
\author{Riposte Laboratories}
\date{}
\ripdocno{RL-Z0-SELF-HOSTING}
\riprev{A}
\ripstrapline{\NoCaseChange{\riphi{The build guide runs master control on a dedicated mini PC. This is the same station, but for people who already have a homelab humming in a closet — a Proxmox node, a TrueNAS box, a rack. Everything here is about \NoCaseChange{\textbf{where the playout stack lives}}; the tower still wants to be a VPS (that's the whole scaling trick), and the storefront lives on Cloudflare. Only the home side changes.}}}
\riprunningtitle{CHANNEL Z0 — Self-Hosting on a Homelab (Proxmox \& TrueNAS)}

\begin{document}

\ripmakehead

\riptoc

\begin{ripnote}
The one hardware fact that decides everything below: ErsatzTV wants a \textbf{hardware H.264 encoder} (Intel Quick Sync / VAAPI, or NVIDIA NVENC). On a homelab that means handing the box's iGPU/GPU to the container. The container stack is \ripmdpath{\ripmono{playout/\allowbreak{}compose.\allowbreak{}yml}}; the rest of this doc is how to give it a GPU on each platform.
\end{ripnote}


\ripsection{\ripmdhead{Where each piece runs}}

\ripcodeopen
\ripcodesize{9.0}{13.95}
\begin{ripcodeblock}
  HOMELAB (Proxmox / TrueNAS)              VPS ($0–8/mo)            CLOUDFLARE
┌───────────────────────────────┐       ┌──────────────┐        ┌──────────────┐
│  playout/compose.yml          │  ONE  │  Owncast     │  HLS   │  Pages        │
│   ├ ersatztv  (needs /dev/dri)│ STREAM│  (the tower) │───────▶│  the          │
│   ├ uplink    (ffmpeg -c copy)│──────▶│  RTMP :1935  │        │  storefront   │
│   └ nowplaying (optional)     │  RTMP │  HLS  :8080  │        │  ch0.ripo…    │
└───────────────────────────────┘       └──────────────┘        └──────────────┘
\end{ripcodeblock}
\ripcodesizereset\ripcodeclose

\textbf{Don't self-host the tower.} It's tempting to run Owncast at home behind a Cloudflare Tunnel and skip the VPS, but: RTMP ingest (:1935) can't traverse Cloudflare's HTTP proxy, and proxying a 24/7 video stream through Pages/Tunnel runs straight into Cloudflare's rules about serving lots of non-HTML traffic. The VPS is \$0 on Oracle's free tier and exists precisely so your house only ever uploads one stream. Keep the split.


\ripsection{\ripmdhead{Proxmox}}

Two ways to host the playout stack. \textbf{Use the LXC path} — for Quick Sync it's dramatically simpler than VM passthrough, and it's lighter.


\ripsubsection{\ripmdhead{Path A — LXC container (recommended)}}

An unprivileged LXC with Docker inside, and the host's \ripmono{/dev/dri} bind-mounted in so the container can reach Quick Sync.

\begin{ripnumbers}
\item \textbf{Create the container.} Debian 12 template, unprivileged, give it 2 cores / 2–4 GB RAM / 20 GB root. Put the media on a separate mount (below).
\item \textbf{Enable Docker-in-LXC + GPU.} Edit \ripmono{/\allowbreak{}etc/\allowbreak{}pve/\allowbreak{}lxc/\allowbreak{}<ID>.\allowbreak{}conf} on the Proxmox host and add: \ripmono{\# Docker needs nesting + keyctl features:\allowbreak{} nesting=\allowbreak{}1,\allowbreak{}keyctl=\allowbreak{}1 \# Hand the iGPU (Quick Sync) to the container lxc.\allowbreak{}cgroup2.\allowbreak{}devices.\allowbreak{}allow:\allowbreak{} c 226:\allowbreak{}* rwm dev0:\allowbreak{} /\allowbreak{}dev/\allowbreak{}dri/\allowbreak{}renderD128,\allowbreak{}gid=\allowbreak{}44} (\ripmono{226} is the DRI major number; \ripmono{gid=44} is usually \ripmono{render} — check with \ripmono{getent group render} on the host and match it. On some hosts the render group is \ripmono{104} or \ripmono{993}.) Restart the container: \ripmono{pct restart <ID>}.
\item \textbf{Media in.} If the library lives on a NAS or another dataset, bind-mount it: \ripmono{pct set <ID> -\allowbreak{}mp0 /\allowbreak{}tank/\allowbreak{}channelz0,\allowbreak{}mp=\allowbreak{}/\allowbreak{}media/\allowbreak{}channelz0}. Or keep it on the container's disk. Either way, point \ripmono{Z0\_MEDIA\_ROOT} at it.
\item \textbf{Install Docker, then the stack.} Inside the container: \ripmono{bash curl -\allowbreak{}fsSL https:\allowbreak{}/\allowbreak{}/\allowbreak{}get.\allowbreak{}docker.\allowbreak{}com | sh git clone <this repo> \&\& cd channel-\allowbreak{}z0/\allowbreak{}playout cp .\allowbreak{}.\allowbreak{}/\allowbreak{}.\allowbreak{}env.\allowbreak{}example .\allowbreak{}env          \# fill in real values docker compose up -\allowbreak{}d             \# ersatztv + uplink}
\item \textbf{Prove Quick Sync reached the container:} \ripmono{bash docker exec -\allowbreak{}it z0-\allowbreak{}ersatztv sh -\allowbreak{}c 'ls -\allowbreak{}l /\allowbreak{}dev/\allowbreak{}dri'   \# renderD128 present? \# optional,\allowbreak{} if vainfo is available in your image:\allowbreak{} \#   vainfo -\allowbreak{}-\allowbreak{}display drm -\allowbreak{}-\allowbreak{}device /\allowbreak{}dev/\allowbreak{}dri/\allowbreak{}renderD128} Then set the ErsatzTV FFmpeg profile to \textbf{VAAPI} and encode a channel; CPU should barely move.
\end{ripnumbers}


\ripsubsection{\ripmdhead{Path B — VM with iGPU passthrough}}

Only worth it if you want full isolation or you're on NVIDIA. Intel iGPU passthrough to a VM means VFIO-binding the iGPU (and the host can't use it for console output), which is fiddly; NVENC in a VM needs the card passed through with the usual \ripmono{hostpci0} + driver dance. If you're already comfortable with GPU passthrough, install Docker in the VM and run \ripmono{playout/\allowbreak{}compose.\allowbreak{}yml} unchanged. If you're not, use Path A — it exists to save you this.


\ripsubsection{\ripmdhead{Proxmox tips}}

\begin{ripbullets}
\item \textbf{Snapshot the ErsatzTV config.} Its whole brain (schedule, filler presets, library index) is the \ripmono{ersatztv-\allowbreak{}config} volume. Put it on a dataset you snapshot; rebuilding the schedule by hand is the one genuinely annoying loss.
\item \textbf{Pin resources.} ErsatzTV idles cheap but spikes at each program transition; 2 cores is plenty, but don't starve it to 1.
\item \textbf{The uplink uses host networking} so it can reach ErsatzTV at \ripmono{127.\allowbreak{}0.\allowbreak{}0.\allowbreak{}1:\allowbreak{}8409}. In an LXC that's fine. If you split them across hosts, set \ripmono{Z0\_\allowbreak{}CHANNEL\_\allowbreak{}URL} to the ErsatzTV LAN address instead.
\end{ripbullets}


\ripsection{\ripmdhead{TrueNAS SCALE}}

TrueNAS SCALE runs Docker apps natively (Dragonfish/Electric Eel and later), and it's a natural home for the station because the \textbf{media library is already on the NAS}. Two moving parts: datasets, and a custom app.


\ripsubsection{\ripmdhead{1 · Datasets}}

Create datasets (not just folders) so you get snapshots and clean ACLs:

\ripcodeopen
\ripcodesize{9.0}{13.95}
\begin{ripcodeblock}
tank/channelz0            # Z0_MEDIA_ROOT  (the whole media tree)
tank/apps/ersatztv        # ErsatzTV config — snapshot this one
\end{ripcodeblock}
\ripcodesizereset\ripcodeclose

Build the media tree inside \ripmono{tank/\allowbreak{}channelz0} with \ripmdpath{\ripmono{tools/\allowbreak{}make-\allowbreak{}media-\allowbreak{}tree.\allowbreak{}sh}} (run it over SSH, or recreate the folders in the UI). Give the \textbf{apps user (uid/gid 568)} read access to the media dataset and read/write to the config dataset — set the ACL in Datasets → Edit Permissions.


\ripsubsection{\ripmdhead{2 · The app}}

Install ErsatzTV + the uplink as a \textbf{Custom App} (Apps → Discover → Custom App, "Install via YAML"), pasting a trimmed \ripmdpath{\ripmono{playout/\allowbreak{}compose.\allowbreak{}yml}} with the dataset paths filled in. The GPU line is the important part — TrueNAS passes \ripmono{/dev/dri} when you declare the device:

\ripcodeopen
\ripcodehead{yaml}
\ripcodesize{9.0}{13.95}
\begin{ripcodeblock}
services:
  ersatztv:
    image: ghcr.io/ersatztv/ersatztv:latest-vaapi
    restart: unless-stopped
    ports: ["8409:8409"]
    volumes:
      - /mnt/tank/apps/ersatztv:/config
      - /mnt/tank/channelz0:/media:ro
    devices:
      - /dev/dri:/dev/dri
    group_add: ["44"]        # the host's 'render' gid, so VAAPI is usable
\end{ripcodeblock}
\ripcodesizereset\ripcodeclose

(Check your render gid with \ripmono{getent group render} in a TrueNAS shell and match \ripmono{group\_add}.) Add the \ripmono{uplink} service from the same compose file, or — if you'd rather keep the relay outside TrueNAS's app manager — run it as a tiny container via Dockge/Portainer, or on the mini PC. The now-playing bridge is optional and can live anywhere that can reach both ErsatzTV and the tower.


\ripsubsection{\ripmdhead{TrueNAS tips}}

\begin{ripbullets}
\item \textbf{Intel only, mostly.} Quick Sync via \ripmono{/dev/dri} is the easy win. For NVIDIA, flip on the NVIDIA drivers in Apps settings and use \ripmono{latest-nvidia}.
\item \textbf{Media on spinning disks is fine} — ErsatzTV reads sequentially. Keep the \textbf{config dataset on an SSD pool} if you have one; the library index likes it.
\item \textbf{Snapshot task} on \ripmono{tank/\allowbreak{}apps/\allowbreak{}ersatztv}, nightly. That's your station's memory.
\item \textbf{UPS}: wire the NAS's UPS into TrueNAS's built-in NUT, and the station rides through power blips like a real broadcaster (build guide, "Reliability").
\item \textbf{Split option}: some people keep TrueNAS as pure storage and run playout on a separate mini PC that mounts \ripmono{tank/\allowbreak{}channelz0} over NFS. Totally valid — the NAS holds the tapes, the little box runs master control. Mount read-only.
\end{ripbullets}


\ripsection{\ripmdhead{Tips \& tricks (any homelab)}}

\begin{ripbullets}
\item \textbf{Verify the encoder before you build a schedule.} A software-encoding fallback will \riphi{work} and then melt a CPU core at 1080p. Confirm VAAPI/NVENC is actually being used in ErsatzTV's transcode log on day one.
\item \textbf{Updates, deliberately.} \ripmono{docker compose pull \&\& docker compose up -\allowbreak{}d} monthly-ish. Watchtower can automate it, but exclude the playout stack from auto-updates before a Ground Zero premiere — never ship on a Friday.
\item \textbf{Two clocks, one truth.} Make sure the playout host's timezone is right; the sign-off ritual and the now-playing bridge both trust the wall clock.
\item \textbf{Health at a glance.} Point UptimeRobot at the tower's \ripmono{https:\allowbreak{}/\allowbreak{}/\allowbreak{}watch.\allowbreak{}<domain>/\allowbreak{}api/\allowbreak{}status}; you'll hear about an outage before a neighbour does.
\item \textbf{Keep the secret in one place.} The stream key lives in \ripmono{.env} (or \ripmono{/\allowbreak{}etc/\allowbreak{}channel-\allowbreak{}z0.\allowbreak{}env}) and nowhere else — not in the compose file, not in a snapshot you export, not in a screenshot of your dashboard.
\end{ripbullets}

\riphi{One stream leaves the house. The homelab just makes the house tidier.}

\ripend

\end{document}
